\clearpage
\fancyhead[L]{\small CARLINI. TAVOLE ASTRONOMICHE.}
\hypertarget{native-548-lower}{}\label{native-548-lower}
\begin{center}\small Göttingische gelehrte Anzeigen. Stück 193. S. 1921..1928. 3. Dezember 1812.\\[3pt]\rule{.72\linewidth}{.3pt}\end{center}

\emph{Esposizione di un nuovo metodo di costruire le tavole astronomiche applicuto alle tavole del sole di} \textsc{Francesco Carlini}. 1810. \emph{Dalla reale stamperia.} LXI und 92 Seiten in Kleinquart. Mailand.

Durch die Arbeiten von \textsc{Laplace}, \textsc{Delambre} und \textsc{von Zach} haben unsere neuesten Sonnentafeln einen solchen Grad an Genauigkeit erreicht, dass in dieser Hinsicht kaum noch etwas zu wünschen bleibt. Bei der Berechnung von Sonnenpositionen, die so oft auszuführen ist und rechnende Astronomen so viel Zeit kostet, wäre aber eine noch \emph{einfachere Handhabung} wünschenswert gewesen. Bei \textsc{von Zach}s Tafeln kam es weniger auf die Handhabung an als darauf, die Tafeln auf kleinem Raum unterzubringen. \textsc{Delambre} hat etwas mehr für die einfache Handhabung getan. Er lieferte die mittleren Bewegungen für alle Tage des Jahres, eine sehr detaillierte Tafel für die Mittelpunktsgleichung und Tafeln mit zwei Eingängen für die Störungen durch die anderen Planeten. Trotzdem lassen diese Tafeln noch viel zu wünschen übrig. Sie wären auch ohne grundlegende Änderungen einfacher zu benutzen gewesen, wenn der Verfasser die Störungen durch den Mond zuvor umgeformt, mehrere kleine Tafeln an ihrem richtigen Platz belassen hätte, statt sie aus sehr unangebrachter Sparsamkeit an anderer Stelle unterzubringen, was jedes Mal zu lästigem Hin- und Herblättern zwingt, und wenn er die Störungen des Radiusvektors in Störungen seines Logarithmus umgeformt hätte. Hinzu kommt, dass die Störungen des Radiusvektors durch Jupiter fehlerhaft und nach einer unvollständigen Formel berechnet sind, sodass man im eigenen Exemplar erst noch eine Korrekturtafel nachtragen muss.

\SourceLeaf{549}Wir freuen uns, den Astronomen mit der vorliegenden Arbeit von \textsc{Carlini} neue Sonnentafeln ankündigen zu können, die in ihrer Handhabung alles übertreffen, was bisher zur Verfügung stand, ja sogar alles, was man sich gewünscht hat. Man kann sagen, dass ihr Verfasser mit ausgefeilter Kunst alles getan hat, um die Benutzung zu erleichtern. Fast jeden Wunsch hat er vorweggenommen. Die Berechnung einer Sonnenposition ist dadurch zu einem leichten Spiel geworden, bei dem man kaum eine Ziffer mehr schreiben muss als unbedingt nötig. Die rechnenden Astronomen schulden dem Verfasser für ein solches Werk den größten Dank. Wir können daher nicht umhin, es etwas ausführlicher zu besprechen.

Der Hauptgedanke, den \textsc{Carlini} so erfolgreich nutzt, besteht darin, die Winkelargumente, von denen eine Tafel abhängt, nach \emph{Tagen} fortschreiten zu lassen, anstelle der sonst üblichen Schritte nach Graden oder Tausendsteln des Umfangs. So lässt sich die Mittelpunktsgleichung der Sonne für eine bestimmte Exzentrizität (\textsc{Carlini} hat die für 1810 gewählt) in einer Tafel angeben, deren Argument die Zahl der Tage seit dem letzten Durchgang durch das Perigäum ist. Da wir noch öfter vom Zeitpunkt eines solchen Durchgangs sprechen müssen, bezeichnen wir ihn der Kürze halber mit $T$. Dieses Argument lässt sich mühelos bilden, wenn man weiß, wie viele Tage zu Beginn jedes Jahres seit dem letzten $T$ vergangen waren oder bis zum nächsten $T$ noch vergehen mussten (\textsc{Carlini} setzt dafür immer den Mittag des unmittelbar vorausgehenden 31. Dezember am Mailänder Meridian an, auch in Schaltjahren). Für jedes Jahr lässt sich dieses Zeitintervall, das im ersten Fall als positiv und im zweiten als negativ gelten kann, als eine Art Epakte in einer eigenen Tafel ansetzen. Es wird durch eine kleine Zahl dargestellt, da $T$ zufällig immer in die Nähe des Jahresanfangs fällt. Damit die Zahl stets positiv ist, hat \textsc{Carlini} sie um zwei Einheiten erhöht und so in Tafel II unter der Überschrift A angesetzt. Also liegen $g$ Tage nach Jahresbeginn $g+A-2$ Tage nach dem nächsten $T$. Die zugehörige Mittelpunktsgleichung, bereits mit der mittleren Bewegung während dieser $g+A-2$ Tage verbunden, steht neben dem Argument $g+A$ in Tafel III unter der Überschrift Prima equazione della longitudine del Sole. Dazu braucht man nur die auf den Anfang des vorgegebenen Jahres bezogene Länge des Perigäums zum Zeitpunkt $T$ zu addieren, um sofort die elliptische Länge der Sonne zu erhalten. Diese Länge steht in Tafel II als Termine constante della longitudine del Sole. Damit wären die beiden zuletzt genannten Rubriken in Tafel III und II vollständig erklärt, wenn die Bewegung der Sonne keine weiteren Ungleichheiten hätte und die Exzentrizität unveränderlich wäre. Da dies nicht der Fall ist, ergeben sich noch einige Änderungen, die wir nun erläutern. Die säkulare Änderung der Exzentrizität macht eine Korrektur der Mittelpunktsgleichung nötig, die sich durch die Formel $(t+f-1810)c$ darstellen lässt. Dabei ist $t$ die Jahreszahl, $f$ der Bruchteil des Jahres und $c$ eine von der mittleren Anomalie abhängige Größe zwischen den Grenzen $+0''173$ und $-0''173$. Setzt man also $c = c^1 - 0''173$, so ist $c^1$ stets positiv. Für die gesamte Korrektur ergibt sich dann \[ = (t-1810)c^1 - (t-1810)0''173 + fc \] Der Teil $-(t-1810)0''173$ ist mit der Länge des Perigäums verbunden und bereits im Termine constante der Tafel II enthalten. Der Teil $fc$ ist mit der Prima equazione der Tafel III verbunden. Daher braucht man nur noch $(t-1810)c^1$ hinzuzufügen. Dazu findet man das stets positive $c^1$ in Tafel IV unter der Überschrift Variazione annua della prima equazione mit dem Argument $g+A$ der dritten Tafel. Mit der Prima equazione in Tafel III ist außerdem der Teil der Nutation verbunden, der von der Länge der Sonne abhängt. Die Tafel selbst geht \SourceLeaf{550}in Schritten von einem Zehnteltag bis 370. Da die erste Differenz $367''$ nicht überschreitet und die zweite ganz verschwindet, lässt sich sehr einfach interpolieren.

Mit der Erklärung der beiden Rubriken der Tafel II sind wir noch nicht ganz fertig. Sehr einfallsreich ist die Art, wie \textsc{Carlini} bei der Berechnung der Aberration deren veränderlichen Teil berücksichtigt. In den bisherigen Tafeln enthält die berechnete Sonnenlänge den konstanten Teil dieser Aberration, $-20''25$. Darin sind also eigentlich \emph{sowohl} die Epoche der mittleren Länge als auch die Länge des Perigäums bereits um $20''25$ vermindert. Sucht man die scheinbare Länge der Sonne, muss man noch den veränderlichen Teil der Aberration aus einer besonderen Tafel hinzufügen. Nach einer Bemerkung von \textsc{Carlini} hätte man sich diesen letzten Schritt erspart, wenn man die Länge des Perigäums nur um die Hälfte, nämlich $10''12$, vermindert und die um $20''25$ verminderte Länge der mittleren Länge beibehalten hätte. Deshalb musste \textsc{Carlini} die Zahlen in der Spalte der Tafel II mit der Überschrift Termine constante um $10''1$ erhöhen und die Zahlen $A$ derselben Tafel um 0,00286 Tag vermindern. So erhält man unmittelbar die Sonnenlänge mit Aberration. Will man dagegen die wahre Sonnenlänge, braucht man nur $10''1$ zum Termine constante und 0,00286 Tag zu $A$ zu addieren. Schließlich enthält der Termine constante noch die Größe $-56''6$, um die Summe aller Konstanten auszugleichen, durch die \textsc{Carlini} die gleich näher betrachteten Störungsgleichungen positiv macht. Diese Summe beträgt $+56''6$.

Die Tafeln für die kleinen Ungleichheiten der Sonnenlänge beruhen auf ähnlichen Überlegungen wie die Tafel für die Mittelpunktsgleichung. Die vom Mond verursachten Ungleichheiten sind durch eine geeignete Umformung auf Gleichungen mit je einem Argument zurückgeführt. Von diesen hat \textsc{Carlini} die beiden größten mit Höchstwerten von $7''5$ und $0''5$ beibehalten. In diesen Tafeln ebenso wie in der Tafel für die Nutation ist das Argument in Tagen angegeben. Diese bezeichnen die Zeit zwischen dem Tag, für den man rechnet, und dem Tag, an dem das Argument der jeweiligen Tafel 0 war. Wie groß jedes so ausgedrückte Argument zu Beginn jedes Jahres von 1750 bis 1900 ist, steht ebenfalls in Tafel II.

Besonders einfach zu benutzen sind schließlich die Tafeln für die Störungen durch die Planeten. \textsc{Carlini} hat sie für zwei Eingänge berechnet, die bekanntlich die mittlere Länge der Erde und die mittlere heliozentrische Elongation des Planeten von der Erde bezeichnen. Beide lassen sich wiederum in Tagen ausdrücken. \textsc{Carlini} hat die Tafel jedoch so eingerichtet, dass man für das erste Argument nur die Zahl der seit Jahresbeginn verstrichenen Tage einzusetzen braucht und für das zweite den Wert, den es zu Jahresbeginn hatte. Dieser ist ebenfalls in Tafel II für die Störungen durch jeden der Planeten Venus, Jupiter, Mars und Saturn angegeben. Das zweite Argument bleibt also das ganze Jahr über gleich. Man entnimmt den Tafeln am besten die Gleichungen für den unmittelbar vorausgehenden und den unmittelbar folgenden Tag, die direkt darin zu finden sind, anstelle der Gleichungen für den Tag selbst, für den man rechnet (damit werden die Tafeln zu Tafeln mit nur einem Eingang), und interpoliert anschließend zwischen den Summen aller vier Störungen. Zu diesen ausgezeichnet eingerichteten Störungstafeln seien uns ein paar kleine Bemerkungen gestattet. \textsc{Carlini} hat sie unter der Voraussetzung berechnet, dass die mittlere Länge der Erde zu Jahresbeginn $= 100^\circ$ sei. Das bedeutet eigentlich, dass man, \emph{wenn man es ganz genau nehmen will,} 1) mit der Zahl der Tage, die seit dem Augenblick vergangen sind, in dem die mittlere Länge der Sonne $= 100^\circ$ war, anstelle der seit Jahresbeginn verstrichenen Tage eingehen muss und 2) mit dem Wert der mittleren Elongation des störenden Planeten genau in diesem Augenblick anstelle des Werts zu Jahresbeginn. Der Unterschied ist allerdings nie nennenswert, \SourceLeaf{551}doch hätte man sich wünschen können, die genau richtigen Argumente jederzeit nach Belieben wählen zu können. Beim ersten Argument wäre dies mit einer besonderen Korrekturtafel möglich gewesen, wie wir sie für den eigenen Gebrauch unserem Exemplar beigefügt haben. Das zweite Argument hätte man in Tafel II ebenso leicht für den genau richtigen Augenblick ansetzen können wie für den Jahresanfang. Im Jahr 1804 ist zum Beispiel die Korrektur des ersten Arguments 1.09, und die genauen Werte von $E$, $F$, $G$, $H$ sind \[ \begin{array}{r|l@{\;}l} E = 36.66 & \text{in der Tafel} & E = 36.55\\ F = 28.45 & & F = 28.34\\ G = 13.13 & & G = 13.09\\ H = 10.1 & & H = 10.1 \end{array} \]

Zweitens hätten wir uns gewünscht, dass die Tafel selbst noch um 10 Tage weiter reicht, also bis 370 Tage anstelle von 360. Wer für die letzten Tage des Dezembers eine Sonnenposition berechnen will, hätte sich so die (freilich geringe) Mühe erspart, die Argumente $E$, $F$, $G$, $H$ zugleich für das folgende Jahr herauszuschreiben.

Um unseren Lesern eine Vorstellung davon zu geben, wie einfach diese Tafeln zu handhaben sind, haben wir den Aufbau der Tafeln für die Längenberechnung etwas ausführlicher beschrieben. Der begrenzte Platz erlaubt es uns nicht, bei den übrigen Tafeln ebenso ausführlich vorzugehen. Wir bemerken daher nur, dass die Tafeln zur unmittelbaren Berechnung des Logarithmus des \emph{Abstands} der Sonne von der Erde und die Tafeln zur Berechnung der \emph{Breite} der Sonne ebenso geschickt angeordnet sind wie die für die Länge. Angefügt sind außerdem einige Tafeln zur Reduktion von der Ekliptik auf den Äquator, für die mittlere Rektaszension der Sonne in Zeit und zur Reduktion der Sonnenpositionen vom mittleren auf den wahren Mittag.

Die Einleitung vor den Tafeln erläutert ihren Aufbau mit Klarheit und Eleganz. Besonders wertvoll ist sie durch die vielen einfallsreichen Kunstgriffe, die sie mitteilt, um die Berechnung einer Ephemeride der Sonnenpositionen für ein ganzes Jahr möglichst zu erleichtern. Der Platz erlaubt es uns hier nicht, auf Einzelheiten einzugehen. Wir halten es aber für unsere Pflicht, allen Astronomen, die viele Sonnenpositionen zu berechnen haben, das Studium und die Nutzung dieser Vorteile nachdrücklich zu empfehlen.
